%%%%%%%%%%%%%%%%%%%%%%%%%%%%%%%%%%%%%%%%%%%%%%%%%%%%%%%
\begin{frame}

\frametitle{Rules}

A Snakemake workflow consists of a Snakemake file ({\texttt{Snakefile}}) that consists of a set of rules.
\vs

Each rule specifies a series of {\blue{directives}}, which includes at least:
\begin{enumerate}
  \item input files,
  \item output files, and
  \item how to generate the output files from the input files
\end{enumerate}
\vs

In many ways, this is similar to Make and Galaxy.

\end{frame}


%%%%%%%%%%%%%%%%%%%%%%%%%%%%%%%%%%%%%%%%%%%%%%%%%%%%%%%
\begin{frame}[fragile]

\frametitle{Sample Rule}

Save the following in {\texttt{Snakefile}}:
\begin{figure}
  \centering
  \begin{tabular}{lc}
%   rule HelloWorld: & \action<+->{{\red{$\leftarrow$}}} \\
  rule HelloWorld: & \only<1>{{\red{$\leftarrow$}}} \\
  {\D\D}input:  "Download/iris.data" & \only<2>{{\red{$\leftarrow$}}} \\
  {\D\D}output:  "HelloWorld/iris.txt" & \only<3>{{\red{$\leftarrow$}}} \\
  {\D\D}shell: & \only<4>{{\red{$\leftarrow$}}} \\
  {\D\D\D\D}""" & \only<5>{{\red{$\leftarrow$}}} \\
  {\D\D\D\D}./helloworld \{input\} {\textgreater}\{output\} & \only<6>{{\red{$\leftarrow$}}} \\
  {\D\D\D\D}"""  & \only<5>{{\red{$\leftarrow$}}} \\
  \end{tabular}
\end{figure}

Note:  Variables are enclosed in curly braces.

\end{frame}



%%%%%%%%%%%%%%%%%%%%%%%%%%%%%%%%%%%%%%%%%%%%%%%%%%%%%%%
\begin{frame}[fragile]

\frametitle{Sample Rule}

{\magenta{\$  }}  {\texttt{snakemake HelloWorld/iris.txt}}
{\tiny{
\begin{figure}
  \centering
  \begin{tabular}{l}
Building DAG of jobs... \\
Provided cores: 1 \\
Rules claiming more threads will be scaled down. \\
Job counts: \\
	count	jobs \\
	1	HelloWorld \\
	1 \\
 \\
rule HelloWorld: \\
    input: Download/iris.data \\
    output: HelloWorld/iris.txt \\
    jobid: 0 \\
 \\
Finished job 0. \\
1 of 1 steps (100\%) done \\
  \end{tabular}
\end{figure}
}}

\end{frame}


%%%%%%%%%%%%%%%%%%%%%%%%%%%%%%%%%%%%%%%%%%%%%%%%%%%%%%%
\begin{frame}[fragile]

\frametitle{Sample Rule}

If the input doesn't exist, then the command will not work:
\vs

{\magenta{\$  }}  {\texttt{snakemake HelloWorld/not-exist.txt}}
\begin{figure}
  \centering
  \begin{tabular}{l}
Building DAG of jobs... \\
MissingInputException in line 1 of Snakefile.py: \\
Missing input files for rule HelloWorld: \\
Download/not-exist.data \\
  \end{tabular}
\end{figure}

Not surprisingly, the {\blue{input}} directive is a precondition for the rule to run.

\end{frame}


%%%%%%%%%%%%%%%%%%%%%%%%%%%%%%%%%%%%%%%%%%%%%%%%%%%%%%%
\begin{frame}

\frametitle{Variable Substitution}

The interesting part of Snakemake (IMHO) occurs during variable substitution; commands can then be generalized.

\end{frame}


%%%%%%%%%%%%%%%%%%%%%%%%%%%%%%%%%%%%%%%%%%%%%%%%%%%%%%%
\begin{frame}

\frametitle{Variable Substitution}

\begin{figure}
  \centering
  \begin{tabular}{lc}
  rule HelloWorld: & \\
  {\D\D}input:  "Download/\{dataset\}.data" & \only<2>{{\red{$\leftarrow$}}} \\
  {\D\D}output:  "HelloWorld/\{dataset\}.txt" & \only<1>{{\red{$\leftarrow$}}} \\
  {\D\D}shell: & \\
  {\D\D\D\D}""" & \\
  {\D\D\D\D}./helloworld \{input\} {\textgreater}\{output\} & \\
  {\D\D\D\D}"""  & \\
  \end{tabular}
\end{figure}


\end{frame}


%%%%%%%%%%%%%%%%%%%%%%%%%%%%%%%%%%%%%%%%%%%%%%%%%%%%%%%
\begin{frame}[fragile]

\frametitle{Variable Substitution}

For example:
\vs

\begin{itemize}
  \item {\magenta{\$  }}  {\texttt{snakemake HelloWorld/abalone.txt}}
  \item {\magenta{\$  }}  {\texttt{snakemake HelloWorld/iris.txt}}
\end{itemize}
\vs

will both work as expected.

\end{frame}


%%%%%%%%%%%%%%%%%%%%%%%%%%%%%%%%%%%%%%%%%%%%%%%%%%%%%%%
\begin{frame}

\frametitle{Top-level Rule}

Instead of typing {\texttt{snakemake}} followed by the output to be generated every time, like Make, there is a top-level rule called {\texttt{all}}.
\vs

It is special in that it has {\red{no}} output defined.

\end{frame}


%%%%%%%%%%%%%%%%%%%%%%%%%%%%%%%%%%%%%%%%%%%%%%%%%%%%%%%
\begin{frame}

\frametitle{Top-level Rule}

\begin{figure}
  \centering
  \begin{tabular}{lc}
  rule all: \\
  {\D\D}input:  "HelloWorld/iris.txt" & \only<3>{{\red{$\leftarrow$}}} \\
  \\
  rule HelloWorld: & \\
  {\D\D}input:  "Download/\{dataset\}.data" & \only<2>{{\red{$\leftarrow$}}} \\
  {\D\D}output:  "HelloWorld/\{dataset\}.txt" & \only<1>{{\red{$\leftarrow$}}} \\
  {\D\D}shell: & \\
  {\D\D\D\D}""" & \\
  {\D\D\D\D}./helloworld \{input\} {\textgreater}\{output\} & \\
  {\D\D\D\D}"""  & \\
  \end{tabular}
\end{figure}

Note:  The {\texttt{all}} rule must come first.

\end{frame}


%%%%%%%%%%%%%%%%%%%%%%%%%%%%%%%%%%%%%%%%%%%%%%%%%%%%%%%
\begin{frame}

\frametitle{Top-level Rule}

{\magenta{\$  }}  {\texttt{snakemake}}
{\tiny{
\begin{figure}
  \centering
  \begin{tabular}{l}
Building DAG of jobs... \\
Provided cores: 1 \\
Rules claiming more threads will be scaled down. \\
Job counts: \\
	count	jobs \\
	1	HelloWorld \\
	1	all \\
	2 \\
 \\
rule HelloWorld: \\
    input: Download/iris.data \\
    output: HelloWorld/iris.txt \\
    jobid: 1 \\
    wildcards: dataset=iris \\
 \\
Finished job 1. \\
1 of 2 steps (50\%) done \\
 \\
localrule all: \\
    input: HelloWorld/iris.txt \\
    jobid: 0 \\
 \\
Finished job 0. \\
2 of 2 steps (100\%) done \\
  \end{tabular}
\end{figure}
}}

\end{frame}


%%%%%%%%%%%%%%%%%%%%%%%%%%%%%%%%%%%%%%%%%%%%%%%%%%%%%%%
\begin{frame}

\frametitle{Multiple Arguments}

\begin{figure}
  \centering
  \begin{tabular}{lc}
  rule HelloWorld: & \\
  {\D\D}input:  "Download/iris.data", \\
  {\D\D\D\D\D\D\D\D}"Download/abalone.data" \\
  {\D\D}output:  "HelloWorld/iris.txt" \\
  {\D\D}shell: \\
  {\D\D\D\D}""" \\
  {\D\D\D\D}\#\#  Assuming helloworld takes two arguments \\
  {\D\D\D\D}./helloworld \{input\} {\textgreater}\{output\} & \only<1>{{\red{$\leftarrow$}}} \\
  {\D\D\D\D}./helloworld \{input[0]\} \{input[1]\} {\textgreater}\{output\} & \only<2>{{\red{$\leftarrow$}}} \\
  {\D\D\D\D}""" \\
  \end{tabular}
\end{figure}

\end{frame}


%%%%%%%%%%%%%%%%%%%%%%%%%%%%%%%%%%%%%%%%%%%%%%%%%%%%%%%
\begin{frame}

\frametitle{Named Inputs/Outputs}

\begin{figure}
  \centering
  \begin{tabular}{lc}
  rule HelloWorld: & \\
  {\D\D}input: \\
  {\D\D\D\D}data1 = "Download/iris.data", \\
  {\D\D\D\D}data2 = "Download/abalone.data" \\
  {\D\D}output:  "HelloWorld/iris.txt" \\
  {\D\D}shell: \\
  {\D\D\D\D}""" \\
  {\D\D\D\D}./helloworld \{input.data2\} {\textgreater}\{output\} & \only<1>{{\red{$\leftarrow$}}} \\
  {\D\D\D\D}""" \\
  \end{tabular}
\end{figure}

\end{frame}


%%%%%%%%%%%%%%%%%%%%%%%%%%%%%%%%%%%%%%%%%%%%%%%%%%%%%%%
\begin{frame}

\frametitle{Parameters}

Within a rule, you might need access to non-file parameters for some reason.  This is where the {\blue{params}} directive is needed:
% {\small{
\begin{figure}
  \centering
  \begin{tabular}{lc}
  rule HelloWorld: \\
  {\D\D}input:  "Download/\{dataset\}.data" \\
  {\D\D}output:  "HelloWorld/\{dataset\}.txt" \\
  {\D\D}params: & \only<1>{{\red{$\leftarrow$}}} \\
  {\D\D}{\D\D}count="42", & \only<2>{{\red{$\leftarrow$}}} \\
  {\D\D}{\D\D}dataset="\{dataset\}" & \only<2>{{\red{$\leftarrow$}}} \\
  {\D\D}shell: \\
  {\D\D}{\D\D}""" \\
  {\D\D}{\D\D}./helloworld --count \{params.count\} --label {\Return} & \only<3>{{\red{$\leftarrow$}}} \\
  {\D\D}{\D\D}{\D\D}\{params.dataset\} \{input\} {\textgreater}\{output\} & \only<3>{{\red{$\leftarrow$}}} \\
  {\D\D}{\D\D}""" \\
  \end{tabular}
\end{figure}

\end{frame}


%%%%%%%%%%%%%%%%%%%%%%%%%%%%%%%%%%%%%%%%%%%%%%%%%%%%%%%
\begin{frame}

\frametitle{Functions for Inputs}

Instead of listing the inputs within a rule, a Python function that returns a list can be used.  ({\red{Finally}}, some Python...  \smiley)
\vs

Functions cannot be defined for outputs, though.

\end{frame}


%%%%%%%%%%%%%%%%%%%%%%%%%%%%%%%%%%%%%%%%%%%%%%%%%%%%%%%
\begin{frame}

\frametitle{Functions for Inputs}

{\small{
\begin{figure}
  \centering
  \begin{tabular}{lc}
  OUTPUT\_DIR = "Download/" & \only<7>{{\red{$\leftarrow$}}} \\
  \\
  def Get\_Inputs (foo): & \only<4>{{\red{$\leftarrow$}}} \\
  {\D\D}my\_inputs = [] \\
  {\D\D}my\_inputs.append (OUTPUT\_DIR + "iris.data") & \only<6>{{\red{$\leftarrow$}}} \\
  {\D\D}my\_inputs.append (OUTPUT\_DIR + "abalone.data") & \only<6>{{\red{$\leftarrow$}}} \\
  {\D\D}print ("Data set:  ", foo.dataset, "{\textbackslash}n") & \only<5>{{\red{$\leftarrow$}}} \\
  {\D\D}return my\_inputs & \only<8>{{\red{$\leftarrow$}}} \\
  \\
  rule HelloWorld: & \only<1>{{\red{$\leftarrow$}}} \\
  {\D\D}input: \\
  {\D\D\D\D}Get\_Inputs & \only<3>{{\red{$\leftarrow$}}} \\
  {\D\D}output:  "HelloWorld/\{dataset\}.txt" & \only<2>{{\red{$\leftarrow$}}} \\
  {\D\D}shell: \\
  {\D\D\D\D}""" \\
  {\D\D\D\D}./helloworld \{input[1]\} {\textgreater}\{output\} & \only<9>{{\red{$\leftarrow$}}} \\
  {\D\D\D\D}""" \\
  \end{tabular}
\end{figure}
}}

\end{frame}


%%%%%%%%%%%%%%%%%%%%%%%%%%%%%%%%%%%%%%%%%%%%%%%%%%%%%%%
\begin{frame}

\frametitle{Functions for Inputs}

Here, we have:
\begin{itemize}
  \item Definition of a simple function, with arguments from the rule passed in.  (Also called {\green{wildcards}} in the documentation.)
  \item Use of a global variable.
\end{itemize}

\end{frame}


%%%%%%%%%%%%%%%%%%%%%%%%%%%%%%%%%%%%%%%%%%%%%%%%%%%%%%%
\begin{frame}

\frametitle{Generating Output with Python}

The {\blue{shell}} directive in a rule states ``how to generate the output files from the input files''.  The shell commands can be replaced with the {\blue{run}} directive followed by Python code.
\vs

Suppose we wanted to count the number of lines in a file:

{\small{
\begin{figure}
  \centering
  \begin{tabular}{lc}
  rule Use\_WC: & \\
  {\D\D}input:  "Download/\{dataset\}.data" & \\
  {\D\D}output:  "Use\_WC/\{dataset\}.txt" & \\
  {\D\D}shell: & \\
  {\D\D}{\D\D}""" & \\
  {\D\D}{\D\D}TMP\_FILE=`mktemp` & \only<1>{{\red{$\leftarrow$}}} \\
  {\D\D}{\D\D}wc -l \{input\} {\textgreater}\{output\} & \\
  {\D\D}{\D\D}rm -f \$\{\{TMP\_FILE\}\} & \only<2>{{\red{$\leftarrow$}}} \\
  {\D\D}{\D\D}""" & \\
  \end{tabular}
\end{figure}
}}

\end{frame}


%%%%%%%%%%%%%%%%%%%%%%%%%%%%%%%%%%%%%%%%%%%%%%%%%%%%%%%
\begin{frame}

\frametitle{Generating Output with Python}

This rule will achieve the same thing:

{\small{
\begin{figure}
  \centering
  \begin{tabular}{lc}
  rule Use\_Python: & \\
  {\D\D}input: & \\
  {\D\D}{\D\D}input\_fn = "Download/\{dataset\}.data"  & \only<2>{{\red{$\leftarrow$}}} \\
  {\D\D}output:  & \\ 
  {\D\D}{\D\D}output\_fn = "Use\_Python/\{dataset\}.txt"  & \\
  {\D\D}run:  & \only<1>{{\red{$\leftarrow$}}} \\
  {\D\D}{\D\D}in\_fn = input.input\_fn  & \only<2>{{\red{$\leftarrow$}}} \\
  {\D\D}{\D\D}in\_fp = open (in\_fn, 'r')  & \\
  {\D\D}{\D\D}line\_count = 0  & \\
  \\
  {\D\D}{\D\D}line = in\_fp.readline ()  & \\
  {\D\D}{\D\D}while line != '':  & \\
  {\D\D}{\D\D}{\D\D}line\_count += 1  & \\
  {\D\D}{\D\D}{\D\D}line = in\_fp.readline ()  & \\
  \end{tabular}
\end{figure}
}}

\end{frame}


%%%%%%%%%%%%%%%%%%%%%%%%%%%%%%%%%%%%%%%%%%%%%%%%%%%%%%%
\begin{frame}

\frametitle{Generating Output with Python}

{\small{
\begin{figure}
  \centering
  \begin{tabular}{lc}
  {\D\D}{\D\D}in\_fp.close ()  & \\
  \\
  {\D\D}{\D\D}\#\#  Write out the value  & \\
  {\D\D}{\D\D}out\_fn = output.output\_fn  & \only<1>{{\red{$\leftarrow$}}} \\
  {\D\D}{\D\D}out\_fp = open (out\_fn, 'w')  & \\
  {\D\D}{\D\D}out\_fp.write (str (line\_count) + " " + in\_fn + "{\textbackslash}n")  & \\
  {\D\D}{\D\D}out\_fp.close ()  & \\
  \end{tabular}
\end{figure}
}}

\end{frame}
